\chapter{GNARL and GNULL: The Two-Layer Runtime}
\label{ch:gnarl-gnull}

When you write \texttt{task}, \texttt{protected}, \texttt{delay until}, or an entry
call in Ada, the compiler does not emit instructions to suspend a thread or to
read a hardware timer. It emits \emph{calls} into the GNAT tasking runtime, and
that runtime is what actually talks to the machine. This chapter is about the
internal architecture of that runtime: two stacked layers named \textbf{GNARL}
and \textbf{GNULL}, and how this bare-metal port wires the bottom of the stack
onto the ESP32-S3 instead of onto an operating system.

The user-facing side --- what the language constructs mean and how you use them
--- is covered in \S\ref{ch:tasking}. Here we look underneath: who is called,
what each layer is responsible for, and which real source units implement it.

\section{Two names, two layers}

The runtime that implements Ada tasking is split into two libraries with two
deliberately retro acronyms.

\begin{itemize}
  \item \textbf{GNARL} --- the \emph{GNat Ada Run-time Library}. This is the
        \emph{upper} layer. It knows Ada semantics: what a task control block is,
        how rendezvous and protected-object queues work, how a task is created,
        activated, and terminated, how priorities and the ceiling protocol
        behave. It is portable: it is written in terms of an abstract notion of
        ``a thread'' and ``the clock,'' and it does not name any specific
        hardware or OS.
  \item \textbf{GNULL} --- the \emph{GNu Low-Level (Library)}. This is the
        \emph{lower} layer. Its one job is to satisfy GNARL's abstract needs ---
        create a thread, find the current thread, set a priority, sleep, wake,
        delay until an absolute time, read a monotonic clock --- by mapping each
        of them onto the services of the underlying platform.
\end{itemize}

The split exists so that GNARL can be written once and ported many times: only
GNULL changes per target. On a hosted GNAT (Linux, Windows) the ``underlying
platform'' GNULL targets is the OS thread library --- pthreads. \textbf{The
single most important idea in this chapter is that here there is no OS, so the
platform GNULL targets is the bare-board kernel} shipped with the bb-runtimes:
the \texttt{System.BB.*} packages. Every place a hosted runtime would call
\texttt{pthread\_create} or \texttt{clock\_gettime}, this port calls into
\texttt{System.BB.Threads} or \texttt{System.BB.Time} instead.

\cnote{On a hosted target the ``OS'' under GNULL is pthreads: \texttt{Self}
becomes \texttt{pthread\_self}, \texttt{Sleep} becomes a condition-variable
wait, the clock is \texttt{clock\_gettime(CLOCK\_MONOTONIC)}. On the ESP32-S3
there is no pthreads and (after FreeRTOS is removed) no scheduler at all but the
one we ship. The bb-runtimes \texttt{System.BB} kernel \emph{is} the OS:
\texttt{System.BB.Threads} provides the ready queue and context switch,
\texttt{System.BB.Time} provides the alarm timer and monotonic clock. GNULL is
the thin Ada shim that names those as ``the OS.''}

\section{The layer cake}

From the source line you write down to the silicon, a tasking operation passes
through five strata. Reading top to bottom:

\begin{enumerate}
  \item \textbf{Ada tasking source} --- \texttt{task}, \texttt{protected},
        \texttt{delay until}, entry calls. The compiler expands these into calls
        into GNARL via \texttt{Rtsfind}.
  \item \textbf{GNARL (upper)} --- \texttt{System.Tasking} and its children. The
        Task Control Block (ATCB), task stages, rendezvous, protected entries.
  \item \textbf{GNULL} --- \texttt{System.Task\_Primitives.Operations} plus the
        abstraction boundary \texttt{System.OS\_Interface}. The primitive set
        (\texttt{Create\_Task}, \texttt{Self}, \texttt{Sleep},
        \texttt{Delay\_Until}, \texttt{Monotonic\_Clock}, \dots).
  \item \textbf{Bare-board kernel} --- \texttt{System.BB.Threads},
        \texttt{System.BB.Time}, \texttt{System.BB.Parameters}, and the Xtensa
        register-level context switch.
  \item \textbf{Hardware} --- the two Xtensa LX7 cores, the system timer that
        drives alarms, and the windowed register file the context switch saves
        and restores.
\end{enumerate}

Table~\ref{tab:gnarl-units} maps the layers to the actual units you will find in
\texttt{crates/esp32s3\_rts/light-tasking-esp32s3/gnarl/}.

\begin{longtable}{@{}p{0.16\linewidth}p{0.34\linewidth}p{0.42\linewidth}@{}}
\caption{The tasking stack, by layer and real source unit}
\label{tab:gnarl-units}\\
\hline
\textbf{Layer} & \textbf{Unit (file)} & \textbf{Role} \\
\hline
\endfirsthead
\hline
\textbf{Layer} & \textbf{Unit (file)} & \textbf{Role} \\
\hline
\endhead
GNARL & \texttt{System.Tasking} (\texttt{s-taskin}) & The ATCB type and task-state model \\
GNARL & \texttt{System.Tasking.Restricted.Stages} (\texttt{s-tarest}) & Task creation/activation under Ravenscar \\
GNARL & \texttt{System.Tasking.Protected\_Objects.Single\_Entry} (\texttt{s-tposen}) & Single-entry protected objects (Ravenscar) \\
GNARL & \texttt{System.Tasking.Stages} (\texttt{s-tassta}), \texttt{...Rendezvous} (\texttt{s-tasren}) & Full-profile task stages and rendezvous (full only) \\
GNULL & \texttt{System.Task\_Primitives.Operations} (\texttt{s-taprop}) & The primitive set GNARL calls \\
GNULL & \texttt{System.OS\_Interface} (\texttt{s-osinte}) & The ``this is the OS'' renaming boundary \\
Kernel & \texttt{System.BB.Threads} (\texttt{s-bbthre}) & Threads, ready queue, scheduling \\
Kernel & \texttt{System.BB.Time} (\texttt{s-bbtime}) & Monotonic clock + alarm queue \\
Kernel & \texttt{System.BB.Parameters} (\texttt{s-bbpara}) & Board constants (\texttt{Ticks\_Per\_Second}, \dots) \\
Kernel & \texttt{context\_switch.S} & Xtensa windowed-register switch (\S\ref{ch:context-switch}) \\
\hline
\end{longtable}

The header comment at the top of \texttt{s-taprop.ads} is the canonical
statement of where the boundary sits. It reads ``GNAT RUN-TIME LIBRARY (GNARL)
COMPONENTS'' and then, in the body of the comment, ``This package contains all
the GNULL primitives that interface directly with the underlying OS.'' That one
file \emph{is} the GNULL spec; everything below it in the table is ``the OS.''

\section{GNARL's upper layer: the ATCB}

The central data structure of GNARL is the \emph{Ada Task Control Block}. In
\texttt{System.Tasking} (\texttt{s-taskin.ads}) it is the type
\texttt{Ada\_Task\_Control\_Block}, and a \texttt{Task\_Id} is simply an access
to one:

\begin{lstlisting}
type Ada_Task_Control_Block;
type Task_Id is access all Ada_Task_Control_Block;
Null_Task : constant Task_Id := null;
\end{lstlisting}

The ATCB holds the Ada-level facts about a task: its state
(\texttt{Unactivated}, \texttt{Runnable}, and so on), its priority, its entry
queues, its activation chain. This is the bare-board ``Jorvik/HI-E version'' of
the package, as its header says --- a restricted, static cut of the full ATCB
with no fields for features the profile forbids. GNARL operates entirely on
\texttt{Task\_Id} values; it never touches a register or a timer directly.
When it needs the machine to actually \emph{do} something --- run this task,
suspend that one --- it calls down into GNULL.

\section{The GNULL primitive set}

\texttt{System.Task\_Primitives.Operations} is a small, fixed menu. Walking the
real spec of the Ravenscar version (\texttt{s-taprop.ads}), the operations GNARL
relies on are:

\begin{lstlisting}
procedure Initialize (Environment_Task : ST.Task_Id);
procedure Create_Task
  (T          : ST.Task_Id;
   Wrapper    : System.Address;
   Stack_Size : System.Parameters.Size_Type;
   Priority   : ST.Extended_Priority;
   Base_CPU   : System.Multiprocessors.CPU_Range;
   Succeeded  : out Boolean);
procedure Enter_Task (Self_ID : ST.Task_Id);
function  Self return ST.Task_Id;
procedure Set_Priority (T : ST.Task_Id; Prio : ST.Extended_Priority);
function  Get_Priority (T : ST.Task_Id) return ST.Extended_Priority;
procedure Sleep (Self_ID : ST.Task_Id; Reason : ST.Task_States);
procedure Wakeup (T : ST.Task_Id; Reason : ST.Task_States);
procedure Delay_Until (Abs_Time : Time);
function  Monotonic_Clock return Time;
function  Is_Task_Context return Boolean;
\end{lstlisting}

Each one is a thin adapter onto the bare kernel. The mapping is the whole point
of the port:

\begin{itemize}
  \item \texttt{Initialize} --- one-time bring-up of the environment task.
        Renames through to \texttt{System.BB.Threads.Initialize}, which sets up
        the ready queue, the alarm timer, and the FPU, and installs the
        environment thread as the first runnable thread.
  \item \texttt{Create\_Task} --- allocate a kernel thread for a new Ada task. It
        calls \texttt{System.BB.Threads.Thread\_Create} with the task's entry
        wrapper, priority, CPU affinity, and a preallocated stack. (Ravenscar
        forbids task termination, so there is deliberately no destroy
        primitive.)
  \item \texttt{Self} --- ``which task is running right now?'' Implemented as
        \texttt{To\_Task\_Id (System.OS\_Interface.Get\_ATCB)} --- the kernel
        stores the ATCB pointer inside each thread descriptor, and
        \texttt{Get\_ATCB} returns the one for the running thread. This is the
        bare-board equivalent of \texttt{pthread\_self} plus thread-local
        storage, but it is a single pointer load.
  \item \texttt{Set\_Priority} / \texttt{Get\_Priority} --- rename onto
        \texttt{System.BB.Threads.Set\_Priority} / \texttt{Get\_Priority}. A
        priority \emph{decrease} (leaving a protected action) is a scheduling
        point; that is where the bare kernel may pick a different thread.
  \item \texttt{Sleep} --- suspend the caller. The body asserts
        \texttt{Self\_ID = Self} (a task can only suspend itself) and then calls
        \texttt{System.OS\_Interface.Sleep}, which renames
        \texttt{System.BB.Threads.Sleep}. The kernel marks the thread
        \texttt{Suspended} and switches away (honoring the \texttt{Wakeup\_Signaled}
        token so a wakeup racing the sleep is not lost).
  \item \texttt{Wakeup} --- make a suspended task runnable again, via
        \texttt{System.BB.Threads.Wakeup}. Note this is \emph{not} itself a
        scheduling point under Ravenscar: the woken task enters the ready queue,
        but the switch happens later when the waker lowers its own priority.
  \item \texttt{Delay\_Until} --- the engine of \texttt{delay until}. It sets the
        caller's state to \texttt{Delay\_Sleep} and arms the kernel alarm queue
        in \texttt{System.BB.Time} for the absolute wake time. When the system
        timer fires at that time, the kernel makes the thread runnable.
  \item \texttt{Monotonic\_Clock} --- the time source behind \texttt{Ada.Real\_Time.Clock}.
        \texttt{System.OS\_Interface} renames \texttt{Clock} onto
        \texttt{System.BB.Time.Clock}, which reads the hardware timer;
        \texttt{Ticks\_Per\_Second} comes from \texttt{System.BB.Parameters}.
  \item \texttt{Is\_Task\_Context} --- ``are we in a task, or in an interrupt
        handler?'' True for task context, False inside an ISR. The kernel tracks
        this with the \texttt{In\_Interrupt} flag on the thread descriptor; the
        restricted runtime uses it to keep tasking operations out of interrupt
        context.
\end{itemize}

The pivot file is \texttt{System.OS\_Interface} (\texttt{s-osinte.ads}). It is
almost entirely \texttt{renames} declarations --- it does no work itself, it just
\emph{relabels} the bare kernel as ``the OS.'' For example:

\begin{lstlisting}
function  Clock return Time renames System.BB.Time.Clock;
procedure Sleep            renames System.BB.Threads.Sleep;
procedure Wakeup (Id : Thread_Id) renames System.BB.Threads.Wakeup;
function  Thread_Self return Thread_Id
                           renames System.BB.Threads.Thread_Self;
function  Get_ATCB return System.Address
                           renames System.BB.Threads.Get_ATCB;
\end{lstlisting}

This is the literal seam between the two worlds: above it, GNARL and GNULL speak
of \texttt{Time} and \texttt{Thread\_Id}; below it, everything is
\texttt{System.BB}. Swapping \texttt{s-osinte} for a pthreads version is, in
principle, all it takes to retarget the runtime to a hosted OS --- which is
exactly how AdaCore ships GNAT on Linux.

\subsection{Where the context switch lives}

The very bottom GNULL primitive --- the moment one task stops running and another
starts --- is not Ada at all. \texttt{Sleep}, \texttt{Delay\_Until}, and priority
changes ultimately reach the kernel's scheduler, which performs the actual
register save/restore in \texttt{context\_switch.S}. On the Xtensa LX7 that means
flushing and restoring the \emph{windowed} register file, not a flat set of
callee-saved registers. That mechanism is intricate enough to have its own
chapter; see \S\ref{ch:context-switch}. From GNULL's point of view the context
switch is simply ``the OS does this''; from the silicon's point of view it is the
single most ESP32-S3-specific line of the whole runtime.

\section{The thread descriptor: GNULL's view of a task}

GNARL has the ATCB; the kernel has its own, leaner record, the
\texttt{Thread\_Descriptor} in \texttt{System.BB.Threads} (\texttt{s-bbthre.ads}).
The two are linked: each \texttt{Thread\_Descriptor} carries an \texttt{ATCB}
field (a \texttt{System.Address}) pointing back at the Ada task, which is what
makes \texttt{Self}/\texttt{Get\_ATCB} a one-pointer lookup. The descriptor holds
exactly what the kernel needs to schedule the thread and nothing more:

\begin{lstlisting}
type Thread_Descriptor is record
   Context         : aliased System.BB.CPU_Primitives.Context_Buffer;
   ATCB            : System.Address;       --  back-pointer to the Ada ATCB
   Base_CPU        : System.Multiprocessors.CPU_Range;
   Base_Priority   : Integer;
   Active_Priority : Integer;              --  raised during protected actions
   Top_Of_Stack    : System.Address;
   Bottom_Of_Stack : System.Address;
   Next            : Thread_Id;            --  ready-queue link
   Alarm_Time      : System.BB.Time.Time;  --  for delay until
   Next_Alarm      : Thread_Id;            --  alarm-queue link
   State           : Thread_States;        --  Runnable | Suspended | Delayed
   In_Interrupt    : Boolean;
   Wakeup_Signaled : Boolean;
   ...
end record;
\end{lstlisting}

Note the deliberate ordering: the comment in the spec insists that
\texttt{Context} sit at the very start of the record, because the assembly
context switch assumes the saved-register buffer is at offset zero of the thread
descriptor. The three-valued \texttt{State} type --- \texttt{Runnable},
\texttt{Suspended}, \texttt{Delayed} --- is the entire scheduling vocabulary of
the restricted profile: a thread is either eligible to run, blocked on an entry,
or waiting out a delay. There is no ``zombie,'' no ``joining,'' no dynamic
priority band, because the profile forbids the features that would need them.

\section{Three profiles, one boundary}

This repository builds three runtime profiles, and the GNARL/GNULL split is the
hinge on which they differ. The selection happens in
\texttt{crates/esp32s3\_rts/gen\_runtime.sh} via the \texttt{ESP32S3\_RTS\_PROFILE}
environment variable; \S\ref{ch:runtime-build} covers the build mechanics in
full.

\begin{itemize}
  \item \textbf{light-tasking} (the default) and \textbf{embedded} both use the
        \emph{restricted} GNARL+GNULL --- the Ravenscar/Jorvik runtime described
        so far. The ATCB is the small \texttt{s-taskin} cut; task creation goes
        through \texttt{s-tarest} (\texttt{Tasking.Restricted.Stages}); protected
        objects are single-entry (\texttt{s-tposen}). There is no rendezvous, no
        \texttt{abort}, no dynamic priority change beyond ceiling locking. The
        two differ only in exception/finalization support
        (\texttt{No\_Exception\_Propagation} + \texttt{No\_Finalization} for
        light-tasking; full ZCX exceptions and finalization for embedded) --- not
        in the tasking model.
  \item \textbf{full} swaps in the \emph{complete} GNARL and a different GNULL.
        \texttt{gen\_runtime.sh} synthesizes it by cloning the embedded source
        tree and overlaying \texttt{crates/esp32s3\_rts/full\_overlay/}: it
        imports the full GNARL units (\texttt{s-tassta}, \texttt{s-tasren} for
        rendezvous, dynamic priorities, abort/ATC machinery) from a native GNAT,
        deletes the restricted stages (\texttt{rm -f s-tarest.*}), and drops in
        \texttt{full\_overlay/gnarl/s-taprop} --- a from-scratch GNULL spec/body
        written against the same \texttt{System.BB} kernel.
\end{itemize}

The full GNULL spec is visibly larger than the Ravenscar one. Its own header
calls it ``the STANDARD Task\_Primitives.Operations interface (as the full GNARL
--- s-tassta / s-tasren --- expects it), implemented over the bare-board BB
kernel.'' Beyond the restricted set it must provide the full lock family
(\texttt{Write\_Lock}/\texttt{Read\_Lock}/\texttt{Unlock} on \texttt{Lock} and
\texttt{RTS\_Lock}), \texttt{Timed\_Sleep}/\texttt{Timed\_Delay},
suspension-object operations, dynamic \texttt{Set\_Priority} with inheritance,
SMP affinity (\texttt{Set\_Task\_Affinity}, \texttt{Get\_CPU}), and the abort/ATC
family (\texttt{Abort\_Task}, \texttt{Suspend\_Task}, \texttt{Stop\_All\_Tasks},
\dots). The crucial design fact: \emph{the bottom of the stack did not change.}
Both the restricted and the full GNULL sit on the same \texttt{System.BB}
kernel. Going to the full profile means a richer upper layer and a fatter
adapter, but the ``OS'' underneath --- threads, clock, alarm queue, context
switch --- is the same bare-board executive in both cases.

\section{Following one operation all the way down}

To make the layering concrete, trace a single \texttt{delay until Next} for a
periodic task:

\begin{enumerate}
  \item \textbf{Source / compiler.} The compiler turns \texttt{delay until Next}
        into a call to a GNARL delay routine, passing the absolute wake time.
  \item \textbf{GNARL.} GNARL converts the \texttt{Ada.Real\_Time.Time} to the
        runtime's internal \texttt{Time} and calls the GNULL primitive
        \texttt{System.Task\_Primitives.Operations.Delay\_Until}.
  \item \textbf{GNULL.} \texttt{Delay\_Until} records
        \texttt{Self\_ID.Common.State := Delay\_Sleep} on the ATCB, then calls
        through \texttt{System.OS\_Interface.Delay\_Until}, which renames
        \texttt{System.BB.Time.Delay\_Until}.
  \item \textbf{Kernel.} \texttt{System.BB.Time} writes the wake time into the
        thread descriptor's \texttt{Alarm\_Time}, links it into the alarm queue
        ordered by expiry, marks the thread \texttt{Delayed}, and programs the
        hardware timer for the nearest alarm. It then invokes the scheduler.
  \item \textbf{Hardware + context switch.} With the caller no longer runnable,
        the scheduler picks the next ready thread and \texttt{context\_switch.S}
        restores its windowed registers. Later, the timer interrupt fires at
        \texttt{Next}, the kernel moves the thread back to \texttt{Runnable}, and
        at the next scheduling point it resumes --- on the line after the
        \texttt{delay until}.
\end{enumerate}

Five layers, one statement. That is the whole reason the runtime is split this
way: GNARL wrote step~2 once, for every Ada target in the world, and this port
only had to supply steps~3--5 by pointing \texttt{System.OS\_Interface} at
\texttt{System.BB} instead of at pthreads. The user view of these same
constructs is in \S\ref{ch:tasking}; the interrupt that ends the delay is in
\S\ref{ch:interrupts}; and the register-level switch in step~5 is in
\S\ref{ch:context-switch}.
